\section{Actual local frames compute the original proper-model line norm}

\begin{definition}[Actual line frame transported through a factor]
\label{def:actual-line-frame-via-factor}
\lean{AlgebraicGeometry.Scheme.Modules.frameFactorIso}
\lean{AlgebraicGeometry.Scheme.Modules.frameViaFactor}
\leanok
For actual morphisms $g:B\to C$, $b:C\to D$, their actual
composite $f$, and an actual unit frame of $b^*M$, transport the
frame to $f^*M$ using the inverse native pullback composition
isomorphism, the actual pullback of the frame by $g$, and the native
pullback-unit isomorphism. Equal composites are transported by the
native pullback equality comparison.
\end{definition}

\begin{lemma}[Actual transported frame composition]
\label{lem:actual-line-frame-via-factor-composition}
\lean{AlgebraicGeometry.Scheme.Modules.frameViaFactor_comp}
\leanok
\uses{def:actual-line-frame-via-factor}
Transporting an actual local frame first through $g$ and then
through $i$ gives the same entire original frame as transporting
through $ig$ at once, with the native original comparison maps.
\end{lemma}
\begin{proof}
The native threefold pullback associativity identity, with its
actual composite equalities, identifies the two whole composition
comparisons. Inverting that equality identifies their inverse
prefixes. Naturality of the actual pullback composition comparison
moves the whole actual original frame across it. The actual
pullback-unit composition identity cancels the remaining unit
comparisons. This proves equality of the actual whole frame
morphisms, hence of their original isomorphisms.
\end{proof}

\begin{corollary}[Original proper-model norm in an actual local frame]
\label{cor:actual-line-model-norm-local-frame}
\lean{EnvelopeProperty.Berkovich.RelativeSchemeValuedPoint.centerOpenFactor}
\lean{EnvelopeProperty.Berkovich.RelativeSchemeValuedPoint.centerOpenFactor_comp}
\lean{EnvelopeProperty.Berkovich.RelativeSchemeValuedPoint.centerLocalFrame}
\lean{EnvelopeProperty.Berkovich.RelativeSchemeValuedPoint.localGenericFactor}
\lean{EnvelopeProperty.Berkovich.RelativeSchemeValuedPoint.localGenericFactor_comp}
\lean{EnvelopeProperty.Berkovich.RelativeSchemeValuedPoint.localFrameGenericIso}
\lean{EnvelopeProperty.Berkovich.RelativeSchemeValuedPoint.centerFrameGenericIso_localFrame}
\lean{EnvelopeProperty.Berkovich.RelativeSchemeValuedPoint.modelLineNorm_eq_localFrame}
\leanok
\uses{lem:actual-local-spectrum-affine-factor,lem:actual-line-frame-via-factor-composition,scheme:proper-native-center,def:scheme-line-actual-model-norm}
An actual unit frame of the original line on any original open
neighborhood of the true proper center computes the norm on the
same original residue-field pullback line. Its coefficient is
computed by pulling that actual local frame along the actual
residue-field morphism into that open.
\end{corollary}
\begin{proof}
The actual proper valuative lift factors through this original
open, because its closed point maps into it. Pull the given
original local frame back along the resulting actual center
factor. Its actual generic frame agrees with the direct original
residue-field pullback frame by the preceding composition
identity and the genuine valuative generic triangle. The original
frame-independence theorem computes the original model norm in
this actual frame. Thus it equals the original residue valuation
of the stated coefficient.
\end{proof}
